% SPDX-License-Identifier: CC-BY-4.0
\section{Formal-verification map and reproduction}\label{sec:verification}

The arguments are complete in the body; reading them requires no Lean. This
appendix locates the formal declarations and gives commands to check them. File
paths are relative to \texttt{formalization/}. The table uses the prefix
\texttt{MR.} for \nolinkurl{MathResearch.}; the external inputs lie under
\nolinkurl{MathResearch.ThirdParty.}

\begingroup
\small
\renewcommand{\arraystretch}{1.3}
\begin{longtable}{@{}>{\raggedright\arraybackslash}p{0.30\textwidth}>{\raggedright\arraybackslash}p{0.65\textwidth}@{}}
\toprule
Paper statement & Lean file and declarations\\
\midrule
\endfirsthead
\toprule
Paper statement & Lean file and declarations\\
\midrule
\endhead
\bottomrule
\endfoot
Hasse multiplicity, Definition~\ref{def:mult}; Lemma~\ref{lem:mult-basic}
& \leanref{claims/HasseMultiplicity.lean}{\nolinkurl{claims/HasseMultiplicity.lean}}\newline
\nolinkurl{MR.HasseMultiplicity.mult}, \nolinkurl{coe_le_mult_iff}, \nolinkurl{min_le_mult_add},
\nolinkurl{add_le_mult_mul}, \nolinkurl{mult_mul}, \nolinkurl{mult_eq_zero_iff}, \nolinkurl{mult_rename_equiv}\\
Digit sums, Lemma~\ref{lem:digits}
& \leanref{claims/TruncatedProduct.lean}{\nolinkurl{claims/TruncatedProduct.lean}}\newline
\nolinkurl{MR.TruncatedProduct.s2_add_le}, \nolinkurl{s2_two_pow}, \nolinkurl{s2_le_self},
\nolinkurl{two_mul_add_s2_le}; \nolinkurl{MR.BinaryMultiplicityDegreeFormula.s2_two_pow_sub_one}\\
Legendre form, Lemma~\ref{lem:legendre}
& \leanref{claims/PerOrderLegendre.lean}{\nolinkurl{claims/PerOrderLegendre.lean}}\newline
\nolinkurl{MR.PerOrderLegendre.legendre}, \nolinkurl{phi_of_lt_two_pow}, \nolinkurl{bits},
\nolinkurl{length_bits}, \nolinkurl{sum_bits}, \nolinkurl{lt_of_mem_bits}\\
Digit-sum identities used in the remarks
& \leanref{claims/BinaryDigitSums.lean}{\nolinkurl{claims/BinaryDigitSums.lean}}\newline
\nolinkurl{MR.BinaryDigitSums.s2_mul_two_pow_add}, \nolinkurl{two_pow_s2_le}, \nolinkurl{s2_le_log},
\nolinkurl{s2_pos}, \nolinkurl{s2_eq_one_iff}, \nolinkurl{s2_succ_add_padicValNat}\\
Catalan parity \cite{AK73,DS06}, Section~\ref{sec:upper}
& \leanref{third-party-claims/CatalanParity.lean}{\nolinkurl{third-party-claims/CatalanParity.lean}}\newline
\nolinkurl{MR.ThirdParty.CatalanParity.padicValNat_catalan}, \nolinkurl{odd_catalan_iff}\\
Telescoping, Lemma~\ref{lem:telescope}; truncated product, Theorem~\ref{thm:truncated-product}
& \leanref{claims/TruncatedProduct.lean}{\nolinkurl{claims/TruncatedProduct.lean}}\newline
\nolinkurl{MR.TruncatedProduct.telescope}, \nolinkurl{g}, \nolinkurl{truncated_product}\\
Hyperplane product, Lemma~\ref{lem:hyperplane-product}; construction, Theorem~\ref{thm:construction}
& \leanref{claims/PerOrderUpperBound.lean}{\nolinkurl{claims/PerOrderUpperBound.lean}}\newline
\nolinkurl{MR.PerOrderUpperBound.hypProd}, \nolinkurl{mult_hypProd}, \nolinkurl{eval_zero_hypProd},
\nolinkurl{totalDegree_hypProd}, \nolinkurl{Phi}, \nolinkurl{construction}, \nolinkurl{construction_spec},
\nolinkurl{per_order_upper_bound}\\
Menezes' polynomial, Proposition~\ref{prop:menezes-polynomial}
& \leanref{claims/PerOrderConstruction.lean}{\nolinkurl{claims/PerOrderConstruction.lean}}\newline
\nolinkurl{MR.PerOrderConstruction.per_order_construction}\\
Expansion, Lemma~\ref{lem:expansion}
& \leanref{claims/PerOrderExpansion.lean}{\nolinkurl{claims/PerOrderExpansion.lean}}\newline
\nolinkurl{MR.PerOrderExpansion.exists_expansion}, \nolinkurl{weight_le_totalDegree}\\
Artin--Schreier root, Section~\ref{sec:lower}
& \leanref{claims/ArtinSchreierSeries.lean}{\nolinkurl{claims/ArtinSchreierSeries.lean}}\newline
\nolinkurl{MR.ArtinSchreierSeries.asRoot}, \nolinkurl{coeff_asRoot_eq_one_iff}, \nolinkurl{asRoot_sq_add},
\nolinkurl{trunc_asRoot}\\
Congruence, Lemma~\ref{lem:congruence}
& \leanref{claims/PerOrderYCoefficient.lean}{\nolinkurl{claims/PerOrderYCoefficient.lean}}\newline
\nolinkurl{MR.PerOrderYCoefficient.xhat}, \nolinkurl{xhat_sq_add}, \nolinkurl{expansion_congr},
\nolinkurl{coeff_expansion}\\
Forms, Lemma~\ref{lem:forms}
& \leanref{claims/PerOrderForms.lean}{\nolinkurl{claims/PerOrderForms.lean}}\newline
\nolinkurl{MR.PerOrderForms.E}, \nolinkurl{E_perm}, \nolinkurl{E_span}, \nolinkurl{E_eval},
\nolinkurl{E_rec}\\
Grid vanishing, Lemma~\ref{lem:grid-vanishing}; grid reduction, Lemma~\ref{lem:grid-reduction}
& Mathlib \nolinkurl{MvPolynomial.eq_zero_of_eval_zero_at_prod_finset};\newline
\leanref{claims/PerOrderSubspaceReduction.lean}{\nolinkurl{claims/PerOrderSubspaceReduction.lean}}\newline
\nolinkurl{MR.PerOrderSubspaceReduction.rho}, \nolinkurl{coeff_reduction}\\
Lowest part, Lemma~\ref{lem:lowest-part}; lower bound, Theorem~\ref{thm:per-order-lower}
& \leanref{claims/PerOrderLowerBound.lean}{\nolinkurl{claims/PerOrderLowerBound.lean}}\newline
\nolinkurl{MR.PerOrderLowerBound.wser}, \nolinkurl{mult_wser}, \nolinkurl{hc_wser_ne},
\nolinkurl{cond}, \nolinkurl{E_vanish}, \nolinkurl{lower_bound}\\
Main theorem, Theorem~\ref{thm:main}
& \leanref{claims/PerOrderValue.lean}{\nolinkurl{claims/PerOrderValue.lean}}\newline
\nolinkurl{MR.PerOrderValue.per_order_value}\\
Minimum over the origin order, Lemma~\ref{lem:minimum}; extremal orders, Remark~\ref{rem:extremal}
& \leanref{claims/PerOrderMinimumDirect.lean}{\nolinkurl{claims/PerOrderMinimumDirect.lean}}\newline
\nolinkurl{MR.PerOrderMinimumDirect.hval}, \nolinkurl{hval_le}, \nolinkurl{hval_attain},
\nolinkurl{min_phi_direct}, \nolinkurl{phi_eq_degreeMin_iff}, \nolinkurl{extremal_of_le_two_pow},
\nolinkurl{extremal_example}, \nolinkurl{delta_eq_D_iff}, \nolinkurl{delta_eq_D_of_le_two_pow},
\nolinkurl{delta_extremal_example};\newline
\leanref{claims/PerOrderMinimum.lean}{\nolinkurl{claims/PerOrderMinimum.lean}}
\nolinkurl{MR.PerOrderMinimum.min_phi}\\
Arithmetic of \(\Phi\): Section~\ref{sec:introduction} (Legendre form, comparison with Menezes'
form, blocks, examples, the top order, the reals), Section~\ref{sec:prelim} (substitutions),
Section~\ref{sec:upper} (exact degree, the factors of \(F\)), Appendix~\ref{sec:second-proof} (one more dimension)
& \leanref{claims/PerOrderArithmetic.lean}{\nolinkurl{claims/PerOrderArithmetic.lean}}\newline
\nolinkurl{MR.PerOrderArithmetic.phi_add_eq}, \nolinkurl{phi_lt_menezes}, \nolinkurl{construction_lt_menezes},
\nolinkurl{phi_add_block}, \nolinkurl{phi_of_le}, \nolinkurl{phi_of_real_range}, \nolinkurl{phi_top},
\nolinkurl{phi_succ}, \nolinkurl{totalDegree_construction}, \nolinkurl{menezes_form_optimal},
\nolinkurl{one_le_s2}, \nolinkurl{example_one}, \nolinkurl{cube_example}, \nolinkurl{example_two},
\nolinkurl{mult_le_mult_zero_subst}, \nolinkurl{card_nonzero}, \nolinkurl{mult_hypProd_sq},
\nolinkurl{totalDegree_hypProd_sq}; for \(\delta\) itself: \nolinkurl{deltaSet}, \nolinkurl{isLeast_deltaSet},
\nolinkurl{delta_succ}, \nolinkurl{delta_of_real_range}, \nolinkurl{delta_top}, \nolinkurl{delta_examples}, \nolinkurl{delta_one}\\
Arithmetic of \(D(n,k)\): Sections~\ref{sec:introduction} and~\ref{sec:main},
Appendix~\ref{sec:second-proof}
& \leanref{claims/DegreeFormulaArithmetic.lean}{\nolinkurl{claims/DegreeFormulaArithmetic.lean}}\newline
\nolinkurl{MR.DegreeFormulaArithmetic.degreeMin_one}, \nolinkurl{degreeMin_four},
\nolinkurl{degreeMin_eight}, \nolinkurl{degreeMin_fifteen}, \nolinkurl{degreeMin_excess},
\nolinkurl{degreeMin_succ}, \nolinkurl{degreeMin_three_mul}, \nolinkurl{degreeMin_add_one_le}; for \(D\)
itself: \nolinkurl{degreeSet}, \nolinkurl{isLeast_degreeSet}, \nolinkurl{admissible_ne_zero},
\nolinkurl{D_examples}, \nolinkurl{D_excess}, \nolinkurl{D_succ}, \nolinkurl{D_three_mul}, \nolinkurl{D_add_one_le}\\
Value of \(D(n,k)\), Corollary~\ref{cor:main}; Theorem~\ref{thm:main-formal}(3)
& \leanref{claims/BinaryMultiplicityDegreeComplete.lean}{\nolinkurl{claims/BinaryMultiplicityDegreeComplete.lean}}\newline
\nolinkurl{MR.BinaryMultiplicityDegreeComplete.degreeMin}, \nolinkurl{degree_complete},
\nolinkurl{cover_bound_attained_iff}\\
Extended range, Corollary~\ref{cor:extended}
& \leanref{claims/BinaryMultiplicityDegreeExtendedRange.lean}{\nolinkurl{claims/BinaryMultiplicityDegreeExtendedRange.lean}}\newline
\nolinkurl{MR.BinaryMultiplicityDegreeExtendedRange.degree_formula_extended}\\
Multiplicity Schwartz--Zippel \cite{DKSS13}, Theorem~\ref{thm:dkss}
& \leanref{third-party-claims/MultiplicitySchwartzZippel.lean}{\nolinkurl{third-party-claims/MultiplicitySchwartzZippel.lean}}\newline
\nolinkurl{MR.ThirdParty.MultiplicitySchwartzZippel.multiplicity_schwartz_zippel}\\
Schwartz--Zippel bound over \(\Ft\), Proposition~\ref{prop:cover-bound}
& \leanref{claims/BinaryMultiplicityCoverBound.lean}{\nolinkurl{claims/BinaryMultiplicityCoverBound.lean}}\newline
\nolinkurl{MR.BinaryMultiplicityCoverBound.cover_bound}, \nolinkurl{cover_bound_of_origin}\\
Linear substitutions, Lemma~\ref{lem:subst}; restriction, Lemma~\ref{lem:restrict}
& \leanref{claims/OneDimensionStep.lean}{\nolinkurl{claims/OneDimensionStep.lean}}\newline
\nolinkurl{MR.OneDimensionStep.homogeneousComponent_subst}, \nolinkurl{totalDegree_subst_le},
\nolinkurl{subst_comp}, \nolinkurl{mult_zero_subst}, \nolinkurl{mult_subst}, \nolinkurl{mult_restrict}\\
Avoiding a linear form, Lemma~\ref{lem:linear-form}
& \leanref{claims/LinearFormAvoidance.lean}{\nolinkurl{claims/LinearFormAvoidance.lean}}\newline
\nolinkurl{MR.LinearFormAvoidance.exists_linForm_not_dvd}, \nolinkurl{exists_subst_linForm_eq_X}\\
One-dimension step, Theorem~\ref{thm:step}
& \leanref{claims/OneDimensionStep.lean}{\nolinkurl{claims/OneDimensionStep.lean}}\newline
\nolinkurl{MR.OneDimensionStep.one_dimension_step}\\
Lower bound, Theorem~\ref{thm:lower}
& \leanref{claims/BinaryMultiplicityLowerBound.lean}{\nolinkurl{claims/BinaryMultiplicityLowerBound.lean}}\newline
\nolinkurl{MR.BinaryMultiplicityLowerBound.lower_bound}\\
Hyperplane covers \cite{BBDM23}, Theorem~\ref{thm:hyperplane}
& \leanref{third-party-claims/HyperplaneCoverUpperBound.lean}{\nolinkurl{third-party-claims/HyperplaneCoverUpperBound.lean}}\newline
\nolinkurl{MR.ThirdParty.HyperplaneCoverUpperBound.hyperplane_cover_upper_bound},
\nolinkurl{card_dot_eq_one}\\
Second proof, Theorem~\ref{thm:main-formal}(1,2); Definition~\ref{def:admissible}
& \leanref{claims/BinaryMultiplicityDegreeFormula.lean}{\nolinkurl{claims/BinaryMultiplicityDegreeFormula.lean}}\newline
\nolinkurl{MR.BinaryMultiplicityDegreeFormula.Admissible}, \nolinkurl{degree_formula},
\nolinkurl{degree_formula_exact}, \nolinkurl{degree_formula_of_lt_two_pow}\\
\end{longtable}
\endgroup

\paragraph{Reading the formal statements.}
\begin{itemize}
 \item Points are functions \(\mathrm{Fin}\,n\to\mathrm{ZMod}\,2\), degree is
 \texttt{MvPolynomial.totalDegree}, and multiplicity is
 \nolinkurl{MR.HasseMultiplicity.mult}, the order of \(P(a+z)\) in
 \(\mathbb N_\infty\). \(\Phi\) is \nolinkurl{MR.PerOrderUpperBound.Phi}, written
 out with natural-number division.
 \item Section~\ref{sec:lower} uses one polynomial ring for both \(x\) and \(y\):
 \(A_\varepsilon(x^2+x)\) is \texttt{aeval yv (A }\(\varepsilon\)\texttt{)}. The
 series \(\hat x_K\) is \texttt{xhat K}, and \(w\) is \texttt{wser k P}, with \(K=k\).
 \item In \nolinkurl{MR.PerOrderForms.E}, the coefficient \(w_{k-1-t-\Sigma}\) is
 written \texttt{hcw k w (t + }\(\Sigma\)\texttt{)}, which is zero when
 \(t+\Sigma\ge k\).
 \item The one-dimension step is stated for the variables
 \(\mathrm{Option}\,\sigma\), with \(\mathrm{none}\) playing the role of \(x_0\).
 The lower bound of Appendix~\ref{sec:second-proof} relabels
 \(\mathrm{Fin}(n+1)\) as \(\mathrm{Option}(\mathrm{Fin}\,n)\).
\end{itemize}

\paragraph{Commands.} The project pins Lean \texttt{4.34.0-rc2} and Mathlib in
\texttt{lakefile.lean} and \texttt{lake-manifest.json}. With
\href{https://lean-lang.org/install/manual/}{elan}, Git and Python~3 installed:

\begingroup\small
\begin{verbatim}
git clone https://github.com/kbr-/math-research.git bmd-check
cd bmd-check
git checkout --detach ab7b3377138bc474e774ba9e4f0f009fd3a74eaa
cd formalization
export MATHLIB_NO_CACHE_ON_UPDATE=1
lake exe cache get claims/BinaryMultiplicityDegreePaper2.lean
python3 verify.py --target claims/BinaryMultiplicityDegreePaper2.lean \
  --out verification.txt
\end{verbatim}
\endgroup

The target is an import-only module whose header lists 65 declarations: the
main ones (\nolinkurl{per_order_value}, \nolinkurl{construction_spec}, both
\nolinkurl{lower_bound} theorems, \nolinkurl{min_phi}, \nolinkurl{min_phi_direct},
\nolinkurl{degree_complete}, \nolinkurl{cover_bound_attained_iff},
\nolinkurl{per_order_construction}) and those behind the remarks and comparisons.
It imports every module in the table, so the build covers the whole chain, and its
axiom report is transitive. The verifier builds the proofs, prints each listed
declaration's type, and rejects unfinished proofs and nonstandard axioms. The
expected axiom report for the main theorem is
\begin{quote}\small\ttfamily
'MathResearch.PerOrderValue.per\_order\_value'\par
depends on axioms: [propext, Classical.choice, Quot.sound]
\end{quote}
and the same for every other listed declaration. These are Lean's standard
foundations.

\paragraph{Kernel replay.} To recheck the compiled proofs, and every imported
declaration including Mathlib's, in a fresh Lean kernel, run
\begingroup\small
\begin{verbatim}
lake env leanchecker --fresh claims.BinaryMultiplicityDegreePaper2
\end{verbatim}
\endgroup
\noindent from \texttt{formalization/}. On 25 September 2026 this passed on the
author's machine in 202 seconds, on sources identical to the pinned commit.

\paragraph{Independent validation.} All builds and replays reported here ran
on the author's machine. No hosted or third-party build of this module has been
recorded; the repository's Lean workflow currently targets a different
preprint.
